\section{Conclusões}\label{cap3}

Responder às seguintes questões na escrita das conclusões:

\begin{itemize}
    
    \item Qual foi o problema em que o trabalho se focou? (2/3 frases)
    
    \item Quais foram as técnicas utilizadas para resolver o problema (4/5 frases)
    
    \item Quais foram os resultados obtidos? (2/3 frases)
    
    \item O que poderia ter sido feito de forma diferente? Ex: Não utilizar o Vader e utilizar uma abordagem baseada em AM para análise de sentimentos? Não utilizar RN "básicas" mas mecanismos mais avançados como transformers? Utilizar uma fonte de dados mais abrangente? (Crítica ao teu trabalho)
    
\end{itemize}


Tratam-se das considerações finais do trabalho, mostrando que os objetivos foram cumpridos e enfatizando as descobertas feitas durante o projeto. Em geral reserva-se um ou dois parágrafos para sugerir trabalhos futuros.

Observe que neste modelo a conclusão é numerada pelo numeral 2, mas o projeto não tem a obrigatoriedade de possuir apenas 2 capítulos. Alias, espera-se que tenha mais que isso.

